\documentclass[11pt]{article}
\usepackage[margin=1in]{geometry}
\usepackage{amsmath,amssymb,amsthm,hyperref}
\newtheorem{lemma}{Lemma}
\newtheorem{theorem}{Theorem}
\newcommand{\Tr}{\operatorname{Tr}}
\newcommand{\nn}[1]{\left\lVert#1\right\rVert}
\newcommand{\uin}[1]{\left|\!\left|\!\left|#1\right|\!\right|\!\right|}
\title{Ando-type product norm inequalities\\MathDB 342339}
\author{Principia Math in Progress}
\date{12 September 2026}
\begin{document}
\maketitle

\section*{Task 1. Resolve both proposed norm inequalities}
Conjecture 2.1 of the published paper \cite{DHLV} (Conjecture 2.4 in
its arXiv version) proposes two inequalities for every unitarily
invariant norm. We prove both in the nonnegative function class used
by its surrounding Theorems 2.1 and 2.2. The conjecture sentence omits
nonnegativity in its operator-convex clause; we also resolve that
literal signed version by an explicit counterexample below.

\begin{theorem}\label{main}
Let $A,B$ be positive semidefinite complex $n\times n$ matrices and
let $f(0)=0$.
If $f$ is operator monotone on $[0,\infty)$, then
\begin{equation}\label{mono}
 \uin{(A-B)(f(A)-f(B))}\leq\uin{|A-B|f(|A-B|)}.
\end{equation}
If $f$ is nonnegative and operator convex on $[0,\infty)$, then
\begin{equation}\label{convex}
 \uin{(A-B)(f(A)-f(B))}\geq\uin{|A-B|f(|A-B|)}.
\end{equation}
Both assertions hold for every unitarily invariant norm.
\end{theorem}

Write $s_1(X)\geq\cdots\geq s_n(X)$ for singular values and
$\nn{X}_{(k)}=\sum_{i=1}^k s_i(X)$ for the Ky Fan norms. By the
Ky Fan dominance theorem, it suffices to prove each assertion for
all $k$. The product $ (A-B)(f(A)-f(B))$ need not be Hermitian.

\subsection*{The operator-monotone assertion}
Put $C=A-B$, $H=f(A)-f(B)$, $x_i=s_i(C)$, and $y_i=s_i(H)$.
Ando's classical norm inequality \cite{Ando} gives
\[
 \sum_{i=1}^j y_i\leq\sum_{i=1}^j f(x_i)\qquad(1\leq j\leq n).
\]
Here $f$ is nonnegative because it is increasing and $f(0)=0$.
Set $z_i=f(x_i)$. The singular-value product inequality and Abel
summation imply
\begin{align*}
 \nn{CH}_{(k)}&\leq\sum_{i=1}^k x_i y_i,\\
 \sum_{i=1}^k x_i(y_i-z_i)
 &=x_k\sum_{i=1}^k(y_i-z_i)
   +\sum_{j=1}^{k-1}(x_j-x_{j+1})\sum_{i=1}^j(y_i-z_i)\leq0.
\end{align*}
Since $x_i f(x_i)$ is decreasing in $i$, this proves
$\nn{CH}_{(k)}\leq\nn{|C|f(|C|)}_{(k)}$, and hence \eqref{mono}.
The product inequality follows by the von Neumann trace inequality
applied to
$\Tr(U^*CHV)=\Tr(HVU^*C)$, where $U,V$ are $n\times k$
isometries: $VU^*C$ has rank at most $k$ and
$s_i(VU^*C)\leq s_i(C)$. Maximizing over $U,V$ gives the claim.

If operator monotonicity is defined to permit a jump at the endpoint
zero, the same argument follows by approximation. In that convention
$f(t)=h(t)+a\mathbf1_{t>0}$, where $a\geq0$ and $h$ is continuous
operator monotone with $h(0)=0$. Apply the continuous case to
$h(t)+at/(t+\varepsilon)$ and let $\varepsilon\downarrow0$ in
finite-dimensional functional calculus.

\subsection*{A weighted resolvent inequality}
The following elementary argument extends the power weights in the
trace method of \cite{DDFL} to spectral cutoffs.

\begin{lemma}\label{frob}
Suppose $D,E$ are commuting positive semidefinite matrices, with
simultaneous diagonal entries $d_i,e_i$ ordered in the same direction.
If $W=W^*$ and $\nn{W}\leq1$, then
\[
 \nn{DWE}_F^2\leq\nn{WDE}_F^2\leq\Tr(D^2E^2),
\]
where $\nn{\cdot}_F$ is the Frobenius norm.
\end{lemma}
\begin{proof}
In a common eigenbasis, $|W_{ij}|^2=|W_{ji}|^2$ and
\[
 \nn{WDE}_F^2-\nn{DWE}_F^2
 =\frac12\sum_{i,j}(d_i^2-d_j^2)(e_i^2-e_j^2)|W_{ij}|^2\geq0.
\]
The other inequality follows from $W^2\leq I$ by taking a trace
against $D^2E^2\geq0$.
\end{proof}

Fix $\rho\geq0$ and $0\leq\theta\leq1$, and define
\begin{equation}\label{cut}
 q(t)=t\bigl(\mathbf1_{t>\rho}+\theta\mathbf1_{t=\rho}\bigr)
 \quad(t\geq0).
\end{equation}
Both $q(t)$ and $q(t)/(t+s)$ are nondecreasing for every $s>0$.

\begin{lemma}\label{ordered}
If $B,C\geq0$, $Q=q(C)$, and $f$ is nonnegative operator convex
with $f(0)=0$, then
\begin{equation}\label{ord}
 \Tr\bigl(Q(f(B+C)-f(B))\bigr)\geq\Tr(Qf(C)).
\end{equation}
\end{lemma}
\begin{proof}
For $s>0$ put
\[
 g_s(t)=\frac{st}{s+t},\qquad
 h_s(t)=\frac{st^2}{s+t}=st-sg_s(t).
\]
Let $X=(C+sI)^{-1/2}$ and introduce the positive contractions
\[
 W=s(B+sI)^{-1},\qquad
 V=(C+sI)^{1/2}(B+C+sI)^{-1}(C+sI)^{1/2}.
\]
The resolvent identity gives
\[
 g_s(B+C)-g_s(B)=sWCXVX.
\]
Using cyclicity of trace, set
\[
 M=XQ^{1/2}WC^{1/2},\qquad N=C^{1/2}XVQ^{1/2}.
\]
Then
\begin{align*}
 \Tr\bigl(Q(g_s(B+C)-g_s(B))\bigr)
 &=s\Tr(MN)\\
 &\leq s\nn{M}_F\nn{N}_F
 \leq s\Tr\bigl(QC(C+sI)^{-1}\bigr)=\Tr(Qg_s(C)).
\end{align*}
To justify the last bound, apply Lemma~\ref{frob} first with
$D=XQ^{1/2}$, $E=C^{1/2}$, and contraction $W$, and then with
$D=C^{1/2}X$, $E=Q^{1/2}$, and contraction $V$.
The two required monotone pairs of squared entries are respectively
$q(t)/(t+s),t$ and $t/(t+s),q(t)$.
Thus \eqref{ord} holds for $h_s=st-sg_s$.
It is equality for $f(t)=t$, and for $f(t)=t^2$ its difference is
\[
 \Tr\bigl(Q((B+C)^2-B^2-C^2)\bigr)=2\Tr(QCB)\geq0.
\]

The standard integral representation of nonnegative operator-convex
functions vanishing at zero is
\begin{equation}\label{repr}
 f(t)=\beta t+\gamma t^2+\int_{(0,\infty)}h_s(t)\,d\mu(s),
 \qquad \beta,\gamma\geq0,\quad \mu\geq0.
\end{equation}
One obtains it by applying the positive operator-monotone
representation to $f(t)/t$; its value at zero is the nonnegative
right derivative $\beta$. This is the representation used in
\cite{DHLV}. Integrating the established inequalities proves
\eqref{ord}. Truncated integrals pass to the full integral by
finite-dimensional functional calculus, since the scalar
representation converges at every eigenvalue.
\end{proof}

\begin{lemma}\label{cross}
If $T,D,Q\geq0$, $QD=0$, and $f$ is as in Lemma~\ref{ordered}, then
\begin{equation}\label{crossineq}
 \Tr\bigl(Q(f(T)-f(T+D))\bigr)\geq0.
\end{equation}
\end{lemma}
\begin{proof}
For the linear and quadratic functions the trace in
\eqref{crossineq} vanishes: every term contains $QD$ or $DQ$
after a cyclic permutation. For $h_s$ it equals
\[
 s\Tr\bigl(Q(g_s(T+D)-g_s(T))\bigr)\geq0.
\]
Indeed, $g_s$ is operator monotone, as follows directly from
inverse order and $g_s(T)=sI-s^2(T+sI)^{-1}$.
Now integrate using \eqref{repr}.
\end{proof}

\subsection*{The operator-convex assertion}
Let $C=A-B=C_+-C_-$, with $C_+C_-=0$, and put
$Z=A+C_-=B+C_+$. For each $k$, choose a cutoff $\rho$ and a
fraction $\theta$ so that
\[
 P=\mathbf1_{|C|>\rho}+\theta\mathbf1_{|C|=\rho},
 \qquad 0\leq P\leq I,\qquad\Tr P=k.
\]
For example, take $\rho$ to be the $k$th largest eigenvalue of
$|C|$ and divide the remaining required trace equally across its
eigenspace. This also covers zero eigenvalues and ties.
With $q$ from \eqref{cut}, put $Q_+=q(C_+)$ and $Q_-=q(C_-)$.
Then $PC=Q_+-Q_-$ and $Q_+C_-=Q_-C_+=0$. Consequently,
\begin{align*}
 \Tr\bigl(PC(f(A)-f(B))\bigr)
 &=\Tr\bigl(Q_+(f(A)-f(Z))\bigr)
   +\Tr\bigl(Q_+(f(Z)-f(B))\bigr)\\
 &\quad+\Tr\bigl(Q_-(f(Z)-f(A))\bigr)
   +\Tr\bigl(Q_-(f(B)-f(Z))\bigr)\\
 &\geq\Tr(Q_+f(C_+))+\Tr(Q_-f(C_-))\\
 &=\Tr\bigl(P|C|f(|C|)\bigr).
\end{align*}
The first and fourth terms are nonnegative by
Lemma~\ref{cross}; the second and third are bounded by
Lemma~\ref{ordered}. Since $f$ is nonnegative convex and
$f(0)=0$, $tf(t)$ is nondecreasing. Thus the last trace equals
$\nn{|C|f(|C|)}_{(k)}$.

For every matrix $Y$, the von Neumann trace inequality gives
$|\Tr(PY)|\leq\nn{Y}_{(k)}$ whenever $0\leq P\leq I$ and
$\Tr P=k$: the singular values of $P$ lie in $[0,1]$ and sum
to $k$. Applying this to $Y=C(f(A)-f(B))$ proves
\[
 \nn{C(f(A)-f(B))}_{(k)}\geq\nn{|C|f(|C|)}_{(k)}.
\]
Ky Fan dominance proves \eqref{convex} and completes the theorem.

\subsection*{The omitted sign assumption}
Without nonnegativity, the operator-convex clause is false already
for scalars. Take $f(t)=t^2-3t$, $A=2$, and $B=1$. This $f$ is
operator convex and $f(0)=0$, but $f(2)=f(1)=-2$, so its proposed
inequality reads $0\geq2$ in the absolute-value norm.
The theorem above supplies the substantive nonnegative version,
as well as settling the literal statement's scope.

\section*{Task 2. Classify equality in the norm inequalities}
Determine equality cases for individual Ky Fan and Schatten norms,
tracking equality in the resolvent contractions and the spectral
cutoff trace bound.

\section*{Task 3. Classify admissible weights}
Determine which scalar weights $q$ preserve the ordered inequality
\eqref{ord}, and whether monotonicity of $q$ and $q(t)/(t+s)$
for all $s>0$ is necessary for a universal statement.

\section*{Task 4. Find the maximal function classes}
Characterize functions satisfying \eqref{mono} or \eqref{convex}
for all positive matrices and all unitarily invariant norms, and
decide how far the operator-monotone and operator-convex
hypotheses can be weakened.

\section*{Task 5. Extend the inequalities to operator ideals}
Establish analogues for compact operators and semifinite
noncommutative integration, specifying the integrability and
approximation assumptions needed for singular-value dominance.

\enlargethispage{2\baselineskip}
\begin{thebibliography}{9}
\bibitem{DHLV} T. H. Dinh, M. T. Ho, C. T. Le, and B. K. Vo,
\emph{Two trace inequalities for operator functions},
Mathematical Inequalities \& Applications 22 (2019), 1021--1026.
\href{https://files.ele-math.com/articles/mia-22-69.pdf}{Published paper},
Conjecture 2.1; \href{https://arxiv.org/abs/1904.01961}{arXiv:1904.01961v1},
Conjecture 2.4.
\bibitem{DDFL} T. H. Dinh, R. Dumitru, J. A. Franco, and C. T. Le,
\emph{On some trace inequalities}, Mathematical Inequalities
\& Applications 23 (2020), 467--476.
\href{https://files.ele-math.com/articles/mia-23-38.pdf}{Published paper}.
\bibitem{Ando} T. Ando, \emph{Comparison of norms
$\uin{f(A)-f(B)}$ and $\uin{f(|A-B|)}$},
Mathematische Zeitschrift 197 (1988), 403--409.
\end{thebibliography}
\end{document}
